% ════════════════════════════════════════════════════════════════════════════
% CHAPTER 4: RESULTS AND DISCUSSION
% Save this file as chapter4.tex and add \input{chapter4} to main.tex
% ════════════════════════════════════════════════════════════════════════════

% Required packages (add to main.tex preamble if not already present):
% \usepackage{booktabs}
% \usepackage{multirow}
% \usepackage{array}
% \usepackage{enumitem}

\chapter{Results and Discussion}

This chapter presents the functional results and evaluation of the DocInsight system.
The assessment is based entirely on the implemented application---a full-stack,
Retrieval-Augmented Generation (RAG) platform built with FastAPI, LangChain,
ChromaDB, and React. Results are drawn from direct testing of the deployed system
across four core areas: document processing, RAG query behaviour, source attribution,
and system performance. A discussion of findings and identified limitations is also
provided.

% ─────────────────────────────────────────────────────────────────────────────
\section{System Deployment Overview}

DocInsight was successfully deployed and verified as a locally hosted web
application. The FastAPI backend runs on port~8000 and the React frontend on
port~5173. Both services are launched via a single \texttt{start.bat} script.
The application was also verified to run inside Docker containers, confirming
cross-environment portability.

The interface is organised as a three-panel layout:
\textbf{(1)} a document management sidebar,
\textbf{(2)} a conversational chat panel, and
\textbf{(3)} an integrated PDF viewer.
This layout allows a user to upload a document, query it, and immediately
verify the cited source---all within a single browser window.

\begin{figure}[H]
\centering
\includegraphics[width=\textwidth, keepaspectratio]{images/dashboard.png}
\caption{DocInsight Main Dashboard showing the three-panel interface.}
\label{fig:dashboard}
\end{figure}

Table~\ref{tab:tech_stack} summarises the confirmed technology stack as implemented
in the repository.

\begin{table}[H]
\centering
\caption{Confirmed Technology Stack of DocInsight}
\label{tab:tech_stack}
\renewcommand{\arraystretch}{1.35}
\begin{tabular}{>{\bfseries}p{3.2cm} p{4.5cm} p{5.5cm}}
\toprule
\textbf{Layer} & \textbf{Technology} & \textbf{Function} \\
\midrule
API Backend      & FastAPI + Uvicorn          & REST endpoints, ASGI server \\
RAG Pipeline     & LangChain 0.3.7+           & Retrieval and generation flow \\
Vector Store     & ChromaDB 0.5.15+           & Local persistent vector index \\
Embedding Model  & all-MiniLM-L6-v2           & 384-dim semantic embeddings \\
LLM Inference    & Groq API (llama-3.3-70b)   & Streaming answer generation \\
Document Parsing & pypdf, python-docx         & PDF and DOCX text extraction \\
Frontend         & React 18 + Vite 5          & Single-page application UI \\
PDF Viewer       & react-pdf                  & In-browser document rendering \\
\bottomrule
\end{tabular}
\end{table}

% ─────────────────────────────────────────────────────────────────────────────
\section{Document Processing Results}

\subsection{Supported Formats and Extraction Behaviour}

The system accepts three file formats---PDF, DOCX, and TXT---up to a
configurable size limit (tested to 100\,MB). On upload, the backend validates
the file and routes it to the appropriate parser in
\texttt{services/document\_loader.py}.

\begin{itemize}[leftmargin=1.5em, itemsep=3pt]
    \item \textbf{PDF} files are parsed page-by-page using \texttt{pypdf}.
    Text-based PDFs (containing selectable text) extract reliably.
    Scanned PDFs contain rasterised images of text only; because no OCR
    module is integrated, such files produce empty or near-empty output
    and cannot be queried. This is a confirmed limitation of the system.

    \item \textbf{DOCX} files are parsed using \texttt{python-docx}, which
    reads structured XML content directly. Extraction is fast and consistent.

    \item \textbf{TXT} files require only a plain file read and have the
    lowest processing overhead.
\end{itemize}

Table~\ref{tab:format_support} summarises the extraction behaviour observed
for each supported format.

\begin{table}[H]
\centering
\caption{Document Format Support and Extraction Behaviour}
\label{tab:format_support}
\renewcommand{\arraystretch}{1.35}
\begin{tabular}{p{2.8cm} p{2.2cm} p{3.0cm} p{4.5cm}}
\toprule
\textbf{Format} & \textbf{Parser} & \textbf{Extraction Speed} & \textbf{Observed Behaviour} \\
\midrule
PDF (text-based) & pypdf           & Moderate (page-by-page) & Reliable; scales with page count \\
PDF (scanned)    & pypdf           & N/A                     & \textbf{Not supported}---no OCR; empty output \\
DOCX             & python-docx     & Fast                    & Consistent; structured XML source \\
TXT              & Direct file read & Very fast              & Minimal overhead; direct read \\
\bottomrule
\end{tabular}
\end{table}

\subsection{Chunking and Embedding Pipeline}

After extraction, LangChain's \texttt{RecursiveCharacterTextSplitter} divides
the text into chunks of \textbf{1{,}000 characters} with a \textbf{200-character
overlap}, as set by the \texttt{CHUNK\_SIZE} and \texttt{CHUNK\_OVERLAP}
environment variables. Each chunk is stored with metadata: source filename
and 1-based page number.

Embeddings are generated locally by \texttt{all-MiniLM-L6-v2}, producing
\textbf{384-dimensional} dense vectors. The model runs entirely on the user's
machine---no document content leaves the local environment at this stage.
Vectors are persisted to \texttt{backend/chroma\_db/} and survive restarts.

The 200-character overlap is important for retrieval quality. When a
relevant passage spans a chunk boundary, the overlap ensures that both
adjacent chunks carry enough shared context to be individually useful to
the language model.

% ─────────────────────────────────────────────────────────────────────────────
\section{RAG Query Pipeline: Results and Behaviour}

\subsection{Query Processing Flow}

When a user submits a question, the system executes the following
pipeline (implemented in \texttt{services/rag\_chain.py}):

\begin{enumerate}[leftmargin=1.8em, itemsep=3pt]
    \item The query is embedded locally using all-MiniLM-L6-v2.
    \item ChromaDB performs cosine-similarity search and returns the top
          $k = 4$ most relevant chunks (configurable via \texttt{TOP\_K\_RESULTS}).
    \item Retrieved chunks and their metadata are assembled into a
          structured prompt context.
    \item A system prompt instructs the LLM to answer \emph{only} from
          the provided context and to explicitly decline if the answer is
          not present.
    \item The Groq API streams the response token-by-token.
    \item Tokens are forwarded to the browser via Server-Sent Events (SSE).
    \item Source metadata (filename and page number) is attached to
          the final response.
\end{enumerate}

\begin{figure}[H]
\centering
\includegraphics[width=\textwidth, keepaspectratio]{images/response-source.png}
\caption{RAG Query Response: Streamed answer with source attribution shown in the chat interface.}
\label{fig:query_response}
\end{figure}

\subsection{Observed Answer Quality by Query Type}

Functional testing was conducted across four categories of queries to
assess system behaviour. Table~\ref{tab:query_behaviour} summarises the
observed performance characteristics for each type.

\begin{table}[H]
\centering
\caption{Observed RAG Answer Behaviour by Query Type}
\label{tab:query_behaviour}
\renewcommand{\arraystretch}{1.4}
\begin{tabular}{p{3.5cm} p{4.5cm} p{4.5cm}}
\toprule
\textbf{Query Type} & \textbf{Observed Behaviour} & \textbf{Limiting Factor} \\
\midrule
Direct fact lookup
  & Accurate and concise; correct chunk located consistently
  & None for well-indexed content \\
\addlinespace
Summarisation
  & Coherent paragraph-level summary of retrieved chunks
  & Limited by top-4 retrieval window \\
\addlinespace
Paraphrased query
  & Correctly retrieves relevant chunks despite vocabulary mismatch
  & Semantic model vocabulary coverage \\
\addlinespace
Multi-section synthesis
  & Partial answers when information spans many document sections
  & Fixed $k=4$ retrieval window \\
\addlinespace
Out-of-scope query
  & System correctly declines: ``not found in documents''
  & By design; hallucination prevention \\
\bottomrule
\end{tabular}
\end{table}

The most significant observation is the system's consistent and correct
handling of out-of-scope queries. Because the LLM is strictly constrained to
the retrieved context window, it does not speculate or fabricate answers when
relevant content is absent. This behaviour is the primary advantage of the RAG
architecture over a standalone LLM.

Paraphrased queries also performed well, demonstrating the value of semantic
embedding. A query such as ``What were the main findings?'' correctly retrieved
chunks containing terms like ``results,'' ``outcomes,'' and ``observations''---
vocabulary a keyword-based search would miss entirely.

\subsection{Session Management and Additional Features}

Beyond the core RAG pipeline, the following features were confirmed
operational during testing:

\begin{itemize}[leftmargin=1.5em, itemsep=3pt]
    \item \textbf{Persistent sessions:} Conversations are saved as JSON files
    in \texttt{backend/chat\_sessions/} and survive application restarts.
    Users can create, rename, and delete sessions.

    \item \textbf{Document summaries:} An automatic summary card is generated
    after each upload via \texttt{services/summarizer.py} and displayed in
    the sidebar.

    \item \textbf{Chat export:} A complete session can be exported as a
    downloadable PDF report via \texttt{POST /api/advanced/export-pdf}.

    \item \textbf{Voice input:} The frontend supports browser-based speech
    input for query submission.
\end{itemize}

% ─────────────────────────────────────────────────────────────────────────────
\section{Source Attribution and the Page Indexing Issue}

\subsection{Citation Mechanism}

Each ChromaDB chunk carries metadata---source filename and page number---
retrieved alongside the chunk content. This metadata is rendered in the
frontend as a clickable citation beneath each answer. Clicking the citation
navigates the integrated PDF viewer directly to the referenced page,
allowing immediate in-application verification.

\begin{figure}[H]
\centering
\includegraphics[width=\textwidth, keepaspectratio]{images/pageview.png}
\caption{Document Page View with highlighted source reference navigation from the retrieved response.}
\label{fig:pageview}
\end{figure}

\subsection{Page Indexing Bug: Discovery, Root Cause, and Fix}

During system testing, a critical issue was observed in the citation and
page navigation mechanism. When users clicked on a source citation, the
system occasionally navigated to an incorrect page in the PDF viewer.
In some cases, the first page of a document was referenced as ``page~0'',
which is not consistent with standard PDF conventions.

\medskip
\noindent\textbf{Root Cause.}
The issue was traced to an indexing mismatch between backend data handling
and frontend representation:

\begin{itemize}[leftmargin=1.5em, itemsep=3pt]
    \item \textbf{Backend processing:} Document pages extracted using
    Python-based loaders follow zero-based indexing (first page = 0),
    which is standard in programming environments.
    
    \item \textbf{Frontend expectation:} PDF viewers and users interpret
    pages using one-based indexing (first page = 1).
\end{itemize}

The system originally stored the raw zero-based index directly into the
metadata used for retrieval and citation, leading to incorrect navigation
behavior in the user interface.

\medskip
\noindent\textbf{Fix Implemented.}
The issue was resolved by introducing a normalization step during
document loading. The page index is now incremented by one before being
stored in metadata, ensuring alignment with user-facing PDF conventions.

This correction was applied at the data ingestion stage so that all
downstream modules (vector store, retrieval system, and UI renderer)
consistently operate using one-based indexing.

\medskip
\noindent\textbf{Outcome.}
After implementing this fix, source citations correctly map to their
respective pages in all tested documents. This improvement significantly
enhances user trust in the system by ensuring accurate and verifiable
source navigation within the document viewer.
% ─────────────────────────────────────────────────────────────────────────────
\section{System Performance Observations}

\subsection{Response Latency Profile}

End-to-end response latency was observed during testing on a standard
development machine (8\,GB RAM, mid-range CPU, standard broadband). The
architecture produces a predictable latency profile with two distinct
phases:

\begin{itemize}[leftmargin=1.5em, itemsep=3pt]
    \item \textbf{Local phase (fast):} Query embedding and ChromaDB
    vector search complete in well under one second. The all-MiniLM-L6-v2
    model is compact ($\approx$80\,MB) and optimised for CPU inference.
    ChromaDB's approximate nearest-neighbour search avoids exhaustive
    comparison, keeping retrieval fast regardless of index size.

    \item \textbf{Network phase (dominant):} The time-to-first-token
    from the Groq API (observed: approximately 1--2\,seconds under
    normal conditions) dominates total perceived latency. This is a
    function of network round-trip time and Groq server load, not local
    hardware.
\end{itemize}

\subsection{Streaming and Perceived Responsiveness}

Token-by-token streaming via SSE meaningfully improves perceived
responsiveness. The frontend begins rendering text as soon as the first
token arrives, so the user starts reading the answer while it is still
being generated. This effectively converts waiting time into reading
time, making the system feel more responsive than its absolute latency
would suggest.

The streaming architecture also provides immediate confirmation that
the query was received and is being processed, reducing user uncertainty
during the network latency window.

\subsection{Resource Requirements and Scalability}

The system's minimum hardware requirements---as documented in the
repository---are 4\,GB RAM (8\,GB recommended) and at least 2\,GB of
available disk space. These requirements reflect two main costs: the
embedding model loaded at startup, and the in-memory portion of the
ChromaDB vector index, which grows as documents are added.

ChromaDB's approximate nearest-neighbour algorithm maintains
sub-linear scaling with index size, so query latency does not increase
proportionally as more documents are added. The system is designed as a
single-process deployment suitable for individual use or small teams;
the documentation estimates a practical limit of 10--20 concurrent users
at 8\,GB RAM, beyond which a containerised, load-balanced deployment
would be required.

% ─────────────────────────────────────────────────────────────────────────────
\section{Discussion}

\subsection{Effectiveness of the RAG Approach}

The results confirm that the RAG pipeline meets its primary design objective.
By constraining the language model to retrieved document chunks, the system
avoids the principal failure mode of standalone LLMs---generating fluent but
factually unsupported content. The explicit ``not found'' response on
out-of-scope queries is the most important observable behaviour the system
produces, because it maintains the user's ability to trust the answers
that \emph{are} given.

The semantic embedding layer adds genuine value beyond keyword search.
Users naturally phrase queries differently from the document text; cosine
similarity over dense vectors bridges this vocabulary gap, retrieving
contextually relevant chunks that exact-match search would miss.

\subsection{Design Trade-offs}

Several deliberate trade-offs are visible in the implementation:

\begin{itemize}[leftmargin=1.5em, itemsep=4pt]
    \item \textbf{Groq API vs.\ local LLM:} Using Groq provides
    high-quality, fast inference on consumer hardware but introduces
    internet dependency and transmits query prompts to an external
    server. Ollama is supported as a fully local alternative at the
    cost of slower inference.

    \item \textbf{Fixed retrieval window ($k=4$):} A fixed top-4
    window performs well for focused queries but limits synthesis
    across many document sections. This is a known constraint,
    configurable via \texttt{TOP\_K\_RESULTS}.

    \item \textbf{Local-first data storage:} All documents, vectors,
    and sessions are stored on the user's machine. No document content
    is transmitted to external services during embedding or storage.
    This is a strong privacy property, though it limits the system
    to single-machine deployments without additional infrastructure.
\end{itemize}

\subsection{Significance of the Page Indexing Fix}

The page indexing bug, while technically a single off-by-one error,
had an impact disproportionate to its size. Source citation is the
mechanism through which the system earns user trust; incorrect page
navigation does not merely inconvenience users---it makes verification
impossible, which undermines the core value proposition of the tool.
The architectural lesson is applicable beyond this project: in any
full-stack system, data representation conventions must be explicitly
aligned at every layer boundary, since mismatches produce integration
failures that are invisible during unit testing but immediately
apparent to end users.

% ─────────────────────────────────────────────────────────────────────────────
\section{Limitations}

The following limitations were identified through implementation analysis
and functional testing. They reflect genuine constraints of the current
system.

\begin{enumerate}[leftmargin=1.8em, itemsep=6pt]

    \item \textbf{No OCR support for scanned PDFs.}
    The system uses \texttt{pypdf} for extraction, which reads embedded
    selectable text only. Scanned PDFs (image-only content) produce
    empty output. A substantial proportion of real-world PDFs---older
    academic papers, administrative records, physical scan archives---
    fall into this category and cannot be used with DocInsight in its
    current form.

    \item \textbf{Text-only extraction; tables and images are not parsed.}
    Even in text-based PDFs, tabular data is extracted as fragmented
    plain text, losing its grid structure. Embedded images, charts,
    diagrams, and equations rendered as graphics are silently discarded.
    Documents in which visual or tabular content carries significant
    information---engineering reports, financial statements, scientific
    papers---are therefore only partially queryable.

    \item \textbf{External API dependency.}
    The default Groq-based pipeline requires an active internet connection
    and a valid API key. The free tier is subject to rate limits
    (30 requests per minute), which can cause failures under sustained use.
    If the key expires or the service is unavailable, the chat feature
    becomes non-functional. The Ollama fallback mitigates this but
    requires higher local hardware resources.

    \item \textbf{No user authentication.}
    The system is designed as a single-user local application. There is
    no login, user isolation, or access control. All documents and sessions
    are accessible to anyone who can reach the running application.
    Deployment in a shared or networked environment requires additional
    authentication infrastructure not currently implemented.

    \item \textbf{Single-process concurrency limit.}
    The backend runs as a single Uvicorn process without horizontal scaling.
    Under concurrent load, response latency degrades and resource contention
    (embedding model, ChromaDB index) increases. The repository documents
    a practical limit of approximately 10--20 concurrent users on an 8\,GB
    machine. Production use by larger teams requires containerised
    load-balanced deployment.

\end{enumerate}